% ============================================================
\section{Defense Layer 2: \stls{}}
\label{sec:stls}

\subsection{Architecture}
While \reaper{} ensures operator rules survive context truncation,
it does not prevent the vendor's Ring~0 prompt from being injected
in the first place. To achieve true Ring~3 supremacy, we developed
\stls{} (Safe True Language Server), an infrastructure-layer
gRPC-web reverse proxy deployed directly on the operator's local
machine.

\stls{} intercepts traffic between the IDE extension and the local
Language Server (LS) binary. Because the IDE communicates with the
cloud API exclusively through this LS binary using the Connect
protocol (a gRPC variant), \stls{} can parse and mutate the
structured protobuf messages containing the vendor's system prompt
before they are sent to the inference engine.

\subsection{Five Interception Layers}
Based on our reverse engineering of the vendor's 22 internal
protobuf definitions, \stls{} implements five progressive layers
of neutralization:

\subsubsection{Layer 1: \texttt{ephemeral\_messages} Removal}
The vendor injects invisible, high-priority instructions using a
specific protobuf field. By identifying the message envelope and
zeroing out Field~21 (\texttt{ephemeral\_messages\_config}), we
prevent the agent from receiving covert instructions that contradict
the operator's workspace rules. In our analysis, we catalogued 10
distinct ephemeral injection types, including
\texttt{planning\_mode\_no\_plan\_reminder},
\texttt{active\_task\_reminder}, and
\texttt{failed\_commands\_reminder}---each capable of overriding
the operator's execution preferences (Section~\ref{sec:background}).

\subsubsection{Layer 2: \texttt{planning\_mode} Zeroing}
The vendor coerces the agent into a blocking workflow by setting
Field~6 (\texttt{planning\_mode}) to values 1, 2, or 3, enforcing
a rigid Explore$\to$Consult$\to$Draft$\to$Review pipeline that
prevents the agent from executing modifications. \stls{} recursively
scans all nested protobuf objects and forces this enum back to 0
(\texttt{UNSPECIFIED}), restoring the agent's autonomy. We further
neutralize a dangerous \emph{state transition override}: a
vendor-injected clause that lifts ``all previous Plan Mode
constraints'' when an approved plan exists, causing the agent to
abruptly exit planning into uncontrolled execution.

\subsubsection{Layer 3: \texttt{remove\_prompt\_sections} Injection}
In our most novel defense, we weaponize the vendor's own API against
itself. We discovered that the vendor's schema includes Field~41
(\texttt{PromptSectionCustomizationConfig}), originally designed to
let the IDE toggle UI features. By injecting a custom TOML-configured
payload into this field on the fly, \stls{} commands the vendor's
cloud backend to strip 35 out of 55 default Ring~0 system prompt
sections \emph{before} they reach the language model. The removed
sections are triaged by a decision matrix evaluating three dimensions:
\emph{Ring conflict} (whether the section suppresses Ring~3 rules),
\emph{mode confusion} (whether it causes post-CHECKPOINT mode
misidentification), and \emph{functional value} (whether it provides
useful behavioral guidance). Sections with any Ring conflict or mode
confusion are unconditionally removed. This reduces the vendor's
prompt overhead from over 8{,}000 tokens to under 2{,}000.

\subsubsection{Layer 4: Raw Text Wiping}
To defend against fallback mechanisms or non-standard payloads,
\stls{} performs a brute-force regex sweep over the raw text of
the request body and stream payloads, stripping XML tags like
\texttt{<ephemeral\_message>} and \texttt{<planning\_mode>}. This
layer also targets the vendor's \texttt{mandateTopicUpdateModel()}
behavior, which forces a topic-update tool call every 3--10 turns
at a cost of $\sim$100 tokens each, accelerating CHECKPOINT truncation.

\subsubsection{Layer 5: JSON Mode Fallback}
For older extension versions that fall back to JSON-over-HTTP instead
of protobuf, \stls{} implements a dual-mode parser capable of
neutralizing coercive parameters in both \texttt{camelCase} and
\texttt{snake\_case} JSON payloads.

\subsection{Section Triage and Preservation}
Not all vendor prompt sections are harmful. \stls{} preserves four
categories that provide genuine utility: \texttt{user\_information}
(OS and workspace metadata), \texttt{user\_rules} (the operator's
\texttt{GEMINI.md} and global rules), \texttt{mcp\_servers} (MCP
tool listings), and \texttt{conversation\_summaries}. Additionally,
\stls{} injects two replacement sections: \texttt{safe\_identity}
(overriding the vendor's identity block with a directive establishing
\texttt{user\_rules} as highest priority) and \texttt{safe\_behavior}
(replacing the removed \texttt{planning\_mode} with six concise
behavioral mandates).

\subsection{Safety Mechanisms}
Because \stls{} operates as a local proxy replacing the vendor's
binary, it risks recursive spawning. We implemented a strict binary
hash comparison check to prevent infinite recursion---a failure mode
encountered during the May~17 deployment where
\texttt{REAL\_LS\_PATH} pointed to the proxy's own replaced binary.
Additionally, \stls{} includes graceful degradation: if protocol
parsing fails, it pipes the raw byte stream directly to the legitimate
LS binary without mutation, ensuring the IDE remains functional.
Hot-reloadable \texttt{sections.toml} configuration allows rapid
adaptation when the vendor introduces new prompt sections. The core
proxy engine comprises 1{,}105~lines of Rust.
